% 图契约：本题的源、边界和推导逐项对应；箭头符号按图注解释。
\documentclass[tikz,border=7pt]{standalone}
\usepackage{fontspec,xeCJK,amsmath,amssymb}
\IfFontExistsTF{Times New Roman}{\setmainfont{Times New Roman}}{\setmainfont{TeX Gyre Termes}}
\IfFontExistsTF{Heiti SC}{\setCJKmainfont[AutoFakeBold=1.5]{Heiti SC}}{\setCJKmainfont[AutoFakeBold=1.5]{FandolHei-Regular}}
\usetikzlibrary{arrows.meta,bending,calc,patterns,positioning}
\definecolor{ink}{HTML}{243745}
\definecolor{blue}{HTML}{275DAB}
\definecolor{green}{HTML}{177C63}
\definecolor{red}{HTML}{BD4935}
\definecolor{orange}{HTML}{AA6C16}
\definecolor{purple}{HTML}{715393}
\tikzset{every node/.style={font=\fontsize{9.3}{12.5}\selectfont,text=ink},
axis/.style={-{Stealth[length=4pt]},line width=.65pt,draw=ink!65},
vec/.style={-{Stealth[length=5pt]},line width=1pt,draw=blue},
short/.style={-{Stealth[length=3pt]},line width=.8pt},
guide/.style={dashed,draw=ink!40,line width=.5pt},
card/.style={rounded corners=2pt,draw=blue!65,fill=blue!4,align=center,text width=6.5cm,inner xsep=9pt,inner ysep=7pt},
title/.style={font=\fontsize{10.5}{14}\selectfont\bfseries,text width=15cm,align=center},
label/.style={font=\fontsize{8.4}{11}\selectfont,align=center},
note/.style={label,text width=6.2cm}}
\begin{document}
\begin{tikzpicture}[x=1cm,y=1cm]
\path[use as bounding box] (-8,-3.5) rectangle(8,3.5);
\node[title] at (0,3.15) {理论图 T31\quad 介质边界条件：统一法向，分别统计电荷};
\begin{scope}[shift={(-4,0)}]\fill[green!6] (-2.4,0) rectangle(2.4,2.1);
\fill[blue!5] (-2.4,-1.9) rectangle(2.4,0);
\draw[ink!65,thick] (-2.4,0)--(2.4,0);
\node[label] at(-1.8,-1.6){介质1 $\varepsilon_1$};\node[label] at(-1.8,1.8){介质2 $\varepsilon_2$};
\draw[vec,orange] (-1.8,-.4)--(-1.8,.7) node[left] {$\mathbf n_{12}$};
\draw[vec,orange] (-.8,0)--(.3,0) node[midway,below] {$\mathbf t$};
\draw[vec] (-.6,-1.6)--(.8,-.5) node[right] {$\mathbf E_1$};
\draw[guide] (-.6,-1.6)--(.8,-1.6)--(.8,-.5);
\draw[vec] (-.6,.3)--(.8,1.3) node[right] {$\mathbf E_2$};
\draw[guide] (-.6,.3)--(.8,.3)--(.8,1.3);
\node[label] at(0,2.45){示意取切向、法向分量均正的情形};
\node[label,red] at(1.6,-.1){$\sigma_f$};
\node[note] at(0,-2.65){法向统一由介质 1 指向介质 2；本图为无理想偶极层的普通界面。};
\end{scope}\node[card] at(3.9,2.05){$\mathbf n_{12}\times(\mathbf E_2-\mathbf E_1)=0$\\切向电场连续};
\node[card] at(3.9,0.7){$\mathbf n_{12}\cdot(\mathbf D_2-\mathbf D_1)=\sigma_f$\\法向电位移只由自由面电荷跃变};
\node[card] at(3.9,-0.65){$\varphi_2=\varphi_1$\\普通单层面电荷不使电势跃变};
\node[card] at(3.9,-2){$\sigma_b=(\mathbf P_1-\mathbf P_2)\cdot\mathbf n_{12}$\\$\varepsilon_0(E_{2n}-E_{1n})=\sigma_f+\sigma_b$};

\end{tikzpicture}\end{document}